\documentclass[11pt]{article}
\usepackage[margin=1in]{geometry}
\usepackage{amsmath,amssymb,amsthm,mathtools}
\usepackage{graphicx}
\usepackage{microtype}
\usepackage[hidelinks]{hyperref}
\newtheorem{theorem}{Theorem}
\newtheorem{proposition}[theorem]{Proposition}
\newtheorem{corollary}[theorem]{Corollary}
\newtheorem{lemma}[theorem]{Lemma}
\newcommand{\cl}{\overline}
\newcommand{\Ann}{\operatorname{Ann}}
\title{BAI-free ideal amenability when the quotient is amenable\\\large A substantial partial answer to arXiv:math/0610137}
\author{Autonomous proof packet, lane 14}
\date{17 August 2026}
\begin{document}
\maketitle

\noindent\textbf{Status.} \texttt{candidate\_partial\_likely\_valid}.  The unrestricted source question is not claimed solved.

\section{Source question}
The source is M.~Eshaghi Gordji, B.~Hayati, and S.~A.~R.~Hosseiniun,
``Derivations into duals of closed ideals of Banach algebras,''
arXiv:math/0610137v2 (2006) \cite{source}.  Their Theorem~5.4 says that if
$I$ is a closed ideal of a Banach algebra $A$, $I$ has a bounded approximate
identity (BAI), and both $I$ and $A/I$ are ideally amenable, then $A$ is
ideally amenable.  On PDF page~13 they ask:

\begin{quote}
\emph{Is the above theorem valid if $I$ does not have any bounded
approximate identity?}
\end{quote}

\begin{center}
\includegraphics[width=\textwidth]{figures/open_problem_crop.png}
\end{center}

\section{Definitions and notation}
If $X$ is a Banach $A$-bimodule, its dual actions are fixed by
\[
 \langle x,a\cdot\phi\rangle=\langle xa,\phi\rangle,
 \qquad
 \langle x,\phi\cdot a\rangle=\langle ax,\phi\rangle.
\]
A derivation $D:A\to X$ satisfies
$D(ab)=a\cdot D(b)+D(a)\cdot b$; it is inner if
$D(a)=a\cdot\phi-\phi\cdot a$ for some $\phi\in X$.
The algebra $A$ is \emph{ideally amenable} if every derivation
$A\to J^*$ is inner for every closed two-sided ideal $J$ of $A$.
It is \emph{amenable} in Johnson's sense if every derivation into the dual
of every Banach $A$-bimodule is inner \cite{johnson}.

For closed ideals $I,J\triangleleft A$, write
\[
 K=I\cap J,
 \qquad E_J=\cl{IJ+JI}^{\,\|\cdot\|}\subseteq K.
\]

\section{Proof intuition}
Ideal amenability of $I$ lets us remove, by an inner perturbation, the part
of a derivation $D:A\to J^*$ visible on $K=I\cap J$.  This initially gives
only that $D(i)$ annihilates $K$ for $i\in I$.  Applying the derivation
identity to products in $I$ upgrades this to $D(I^2)=0$, and weak
amenability of $I$ makes $I^2$ dense in $I$.

The adjusted derivation therefore descends to $A/I$.  Its exact coefficient
module is not usually the dual of a quotient ideal: it is
$(J/E_J)^*$, where $E_J=\cl{IJ+JI}$.  This explains the open problem.
Ordinary amenability of $A/I$ controls every dual module and immediately
finishes the argument.  Alternatively, if $E_J=I\cap J$ and $I+J$ is
closed, then $J/E_J\cong (I+J)/I$ is a closed ideal of $A/I$, so ideal
amenability suffices.  The BAI in the source theorem supplies exactly these
two structural properties.  What remains open is whether endpoint ideal
amenability supplies some different way around the residual module.

\section{The BAI-free theorem}
\begin{theorem}[Amenable quotient]\label{thm:amenable-quotient}
Let $I$ be a closed two-sided ideal of a Banach algebra $A$.  If $I$ is
ideally amenable and $A/I$ is amenable, then $A$ is ideally amenable.  No
approximate identity is assumed on $I$.
\end{theorem}

\begin{proof}
Let $J\triangleleft A$ be closed and let $D:A\to J^*$ be a continuous
derivation.  Put $K=I\cap J$, and let
$r:J^*\to K^*$ be restriction.  Since $K$ is a closed ideal of $I$ and $r$
is an $I$-bimodule morphism, $rD|_I:I\to K^*$ is a derivation.  Ideal
amenability of $I$ gives $\lambda\in K^*$ such that
\[
 rD(i)=i\cdot\lambda-\lambda\cdot i \qquad(i\in I).
\]
Choose a Hahn--Banach extension $\Lambda\in J^*$ of $\lambda$, and set
$D_0=D-\delta_\Lambda$.  Then
\begin{equation}\label{eq:restricted}
 D_0(i)|_K=0 \qquad(i\in I).
\end{equation}

For $a,b\in I$ and $j\in J$, the dual actions and the derivation identity
give
\[
 \langle j,D_0(ab)\rangle
 =\langle ja,D_0(b)\rangle+\langle bj,D_0(a)\rangle=0,
\]
because $ja,bj\in I\cap J=K$.  Hence $D_0(I^2)=0$.  Ideal amenability of
$I$ implies weak amenability.  Weak amenability implies
$\cl{I^2}=I$: indeed, if $0\ne f\in I^*$ annihilated $I^2$, then
$T(x)=f(x)f$ would be a nonzero derivation $I\to I^*$; it cannot be inner,
because the bilinear form $\langle y,T(x)\rangle=f(x)f(y)$ is symmetric,
whereas the form associated to an inner derivation is antisymmetric.
Thus continuity yields
\begin{equation}\label{eq:vanish-I}
 D_0|_I=0.
\end{equation}

Let $E=E_J=\cl{IJ+JI}$.  For $i\in I$ and $a\in A$,
\[
 0=D_0(ia)=i\cdot D_0(a),
 \qquad
 0=D_0(ai)=D_0(a)\cdot i.
\]
Pairing with elements of $J$ shows that $D_0(a)$ annihilates $JI$ and
$IJ$, respectively.  Therefore
\[
 D_0(A)\subseteq E^\perp\cong (J/E)^*.
\]
The closed space $E$ is an $A$-submodule of $J$, and $I$ acts trivially on
$J/E$.  Consequently $J/E$ is a Banach $(A/I)$-bimodule.  By
\eqref{eq:vanish-I},
\[
 \bar D:A/I\longrightarrow (J/E)^*,\qquad
 \bar D(a+I)=D_0(a),
\]
is a well-defined continuous derivation.  Amenability of $A/I$ makes
$\bar D$ inner, implemented by some $\mu\in (J/E)^*=E^\perp\subseteq J^*$.
Thus $D_0=\delta_\mu$, and
$D=\delta_{\Lambda+\mu}$ is inner.  Since $J$ and $D$ were arbitrary, $A$
is ideally amenable.
\end{proof}

\section{The exact structural criterion}
\begin{theorem}[Closed-sum/essential-intersection criterion]\label{thm:criterion}
Let $I\triangleleft A$ be closed.  Suppose $I$ and $A/I$ are ideally
amenable and, for every closed ideal $J\triangleleft A$,
\begin{equation}\label{eq:structural}
 I+J\text{ is closed},
 \qquad
 I\cap J=\cl{IJ+JI}.
\end{equation}
Then $A$ is ideally amenable.
\end{theorem}

\begin{proof}
Repeat the proof of Theorem~\ref{thm:amenable-quotient} through the
factorization
\[
 \bar D:A/I\longrightarrow (J/E_J)^*.
\]
By \eqref{eq:structural}, $E_J=I\cap J$.  The natural map
\[
 \Theta:J/(I\cap J)\longrightarrow (I+J)/I,
 \qquad j+(I\cap J)\longmapsto j+I,
\]
is a continuous bijective bimodule map.  Its domain is Banach, and its
range is a closed ideal of $A/I$ because $I+J$ is closed.  The open mapping
theorem makes $\Theta$ a topological bimodule isomorphism.  Hence
$(J/E_J)^*$ is isomorphic, as an $(A/I)$-bimodule, to the dual of the closed
ideal $(I+J)/I$.  Ideal amenability of $A/I$ makes $\bar D$ inner.  The last
step of Theorem~\ref{thm:amenable-quotient} completes the proof.
\end{proof}

\begin{proposition}[Recovery of the source theorem]\label{prop:bai}
If $I$ has a bounded two-sided approximate identity, then
\eqref{eq:structural} holds for every closed ideal $J\triangleleft A$.
\end{proposition}

\begin{proof}
Let $(e_\alpha)$ be a BAI for $I$ with $\sup_\alpha\|e_\alpha\|\le M$.
For $k\in I\cap J$, $e_\alpha k\to k$ and $e_\alpha k\in IJ$; hence
$I\cap J=\cl{IJ}\subseteq\cl{IJ+JI}$, while the reverse inclusion is
automatic.

To prove $I+J$ closed, first note that for $j\in J$, $i\in I$, and
$\varepsilon>0$, one may choose $\alpha$ so that
$\|i-e_\alpha i\|<\varepsilon$.  Since $e_\alpha j\in I\cap J$,
\[
 \operatorname{dist}(j,I\cap J)
 \le \|j-e_\alpha j\|
 \le (1+M)\|j-i\|+\varepsilon.
\]
Taking the infimum over $i$ and then $\varepsilon\downarrow0$ shows that
the injective map
$J/(I\cap J)\to A/I$ is bounded below.  Its range $(I+J)/I$ is therefore
closed.  Thus $I+J$ is closed.
\end{proof}

\begin{corollary}\label{cor:simple}
Suppose $I$ is topologically simple and ideally amenable, $A/I$ is ideally
amenable, and the image in $A/I$ of every closed ideal of $A$ is closed
(in particular, this holds when $A/I$ is finite-dimensional).  Then $A$ is
ideally amenable, without a BAI assumption on $I$.
\end{corollary}

\begin{proof}
For each closed $J\triangleleft A$, topological simplicity gives
$I\cap J=0$ or $I$.  In the first case $E_J=0$.  In the second,
$\cl{I^2}=I$ by ideal amenability, so $E_J=I$.  Thus
$E_J=I\cap J$.  Closedness of the image $(I+J)/I$ implies closedness of
$I+J$, and Theorem~\ref{thm:criterion} applies.
\end{proof}

\section{A genuine non-BAI example}
Let $H$ be an infinite-dimensional Hilbert space and let
$\mathcal S_1(H)$ be the trace-class operators with composition and trace
norm.  This algebra is topologically simple: if a nonzero closed ideal
contains $T\ne0$, rank-one operators on the left and right sandwich $T$ to
produce every rank-one operator, and finite-rank operators are trace-norm
dense.

It is also weakly amenable.  Given a derivation
$D:\mathcal S_1(H)\to\mathcal S_1(H)^*=\mathcal B(H)$, fix a unit vector
$e$ and write $\theta_{x,e}z=\langle z,e\rangle x$.  The formula
\[
 Sx=-D(\theta_{x,e})e
\]
defines a bounded $S\in\mathcal B(H)$.  From
$a\theta_{x,e}=\theta_{ax,e}$ and the derivation identity, evaluated at
$e$, one obtains
\[
 D(a)x=aSx-Sax \qquad(a\in\mathcal S_1(H),\ x\in H).
\]
Thus $D(a)=aS-Sa$ is inner.  Topological simplicity now makes
$\mathcal S_1(H)$ ideally amenable, since its only closed ideals are $0$
and itself.

On the other hand $\mathcal S_1(H)$ has no bounded approximate identity.
If $(u_\alpha)$ were one with $\|u_\alpha\|_1\le C$, choose orthonormal
$e_1,\dots,e_N$ and the rank-one projections $p_k$ onto $\mathbb Ce_k$.
For one $\alpha$ sufficiently far out,
$\|u_\alpha p_k-p_k\|_1<1/2$ for all $k\le N$, so
\[
 C\ge\|u_\alpha\|_1
 \ge \sum_{k=1}^N |\langle u_\alpha e_k,e_k\rangle|
 >N/2,
\]
a contradiction for large $N$.  Therefore Theorem~\ref{thm:amenable-quotient}
applies non-vacuously to the unitization
$\mathcal S_1(H)^\#$, whose quotient by $\mathcal S_1(H)$ is $\mathbb C$.

\section{Why this does not settle the full question}
The common reduction produces $(J/E_J)^*$, not automatically the dual of a
closed ideal of $A/I$.  Endpoint ideal amenability alone is therefore not
enough for the proof given here.  Eight focused upgrades tested whether
$E_J=I\cap J$ or closedness of $I+J$ follows cohomologically, audited a
later claimed full theorem, and explored module-extension counterexamples.
No proof of either missing relation and no counterexample was found.

Ranjbari and Rejali's Theorem~4.1 \cite{ranjbari-rejali} states a stronger
Fr\'echet-algebra theorem that would specialize to the full Banach result.
Its proof first treats an element of $(I\cap J)^*$ as evaluable on arbitrary
elements of $J$.  Repairing that step gives the valid conclusion
$D_0(I^2)=0$ used above, but the final appeal to the weak-amenability
three-space proof still does not identify $(J/E_J)^*$ as an ideal dual.
The exact audit is recorded in the companion proof-gap packet.  Accordingly,
the unrestricted question remains open in this investigation.

\section{Verification and novelty}
The source statement, all module actions, the density-of-products step, both
annihilator calculations, and the quotient-module identification were
audited explicitly.  The BAI recovery and trace-class example are proved in
the packet rather than inferred from analogy.  No computation is used.

The run indexes, exact-title searches, close keyword searches, and the later
2017 primary PDF were checked.  No inspected source stated the two partial
theorems above with a valid proof.  Novelty confidence is moderate; exhaustive
priority is not claimed.  Human review should focus on
\eqref{eq:restricted}--\eqref{eq:vanish-I} and on the topological
identification in Theorem~\ref{thm:criterion}.

\begin{thebibliography}{9}
\bibitem{source}
M. Eshaghi Gordji, B. Hayati, and S. A. R. Hosseiniun,
\emph{Derivations into duals of closed ideals of Banach algebras},
arXiv:math/0610137v2 (2006).

\bibitem{johnson}
B. E. Johnson,
\emph{Cohomology in Banach algebras},
Memoirs of the American Mathematical Society 127 (1972), 1--96.

\bibitem{ranjbari-rejali}
A. Ranjbari and A. Rejali,
\emph{Ideal Amenability of Fr\'echet Algebras},
U.P.B. Scientific Bulletin, Series A 79(4) (2017), 51--60.
\end{thebibliography}

\end{document}
